\begingroup
\makeatletter\tagsleft@true\let\veqno\@@leqno\makeatother
\clearpage\SourceLeaf{326b}\ArithmeticHead{319}{R}{{\scriptsize ABBILDUNG DER ELLIPSE AUF DEN KREIS.}}
\NumberedItem{8}
\noindent {\small[S. 226]}
\par\indent Das berichtigte Resultat[*)] in einfachster Form ist, $\mu\mu=\lambda$ [gesetzt,]
\begin{equation}
\left\{\begin{aligned}
1+2\mu^4+2\mu^{16}+2\mu^{36}+{.}{.}{.}&=a\\
2\mu+2\mu^9+2\mu^{25}+{.}{.}{.}&=b
\end{aligned}\right.\tag*{[19a]}
\end{equation}
\begin{equation}
\int\frac{dt}{\surd(aa-bbtt)(bb-aatt)}=\int\frac{dT}{\surd(4\mu\mu-TT)}\ \text{[**)]},\tag*{[18a]}
\end{equation}
\noindent also
\begin{equation}
\left\{\begin{aligned}
&t+\frac16\left(\frac{bb}{aa}+\frac{aa}{bb}\right)t^3+\left(\frac3{40}\frac{b^4}{a^4}+\frac1{20}+\frac3{40}\frac{a^4}{b^4}\right)t^5+{.}{.}{.}\\[4pt]
&\qquad=\frac{ab}{2\mu}\left\{T+\frac1{24\mu\mu}T^3+\frac3{640\mu^4}T^5+{.}{.}{.}\right\},
\end{aligned}\right.\tag*{[26]}
\end{equation}
\noindent also erste Glieder
\begin{equation}
T=\frac{2\mu}{ab}\left\{t+\frac{a^4+b^4-1}{6aabb}t^3+{.}{.}{.}\right\}.\tag*{[27]}
\end{equation}
\par\bigskip\noindent\rule{.25\linewidth}{.4pt}\par\smallskip
{\small\par\indent [*) Siehe die Gleichungen [18], [19], [20], oben S. 317, 318.]\par
\indent [**) In der Handschrift fehlt rechts vom Gleichheitszeichen das Integralzeichen.]\par}
\clearpage\SourceLeaf{327}\ArithmeticHead{320}{L}{ANALYSIS.\quad NACHLASS.}
\par\indent Auch wird aus dem S. 223 [der Handschrift] gegebenen Product-Ausdruck [Gleichung [15], oben S. 317], wenn man zuerst
\[u=\mu x\]
\noindent schreibt, woraus
\[x+\frac1x=\frac T\mu\]
\noindent wird, leicht abgeleitet
\[t=T\frac{(1-\mu^4)^2+\mu^6TT}{(1-\mu\mu)^2+\mu\mu TT}.\frac{(1-\mu^8)^2+\mu^{14}TT}{(1-\mu^6)^2+\mu^{10}TT}.\text{etc.}\]
\noindent oder
{\small\[\begin{aligned}
t={}&\frac{(1-\mu^4)(1-\mu^8)(1-\mu^{12})\,{.}{.}{.}}{(1-\mu\mu)(1-\mu^6)(1-\mu^{10})\,{.}{.}{.}}.T.\\[6pt]
&\quad\frac{1+\dfrac{\mu^6}{(1-\mu^4)^2}TT}{1+\dfrac{\mu\mu}{(1-\mu\mu)^2}TT}.\frac{1+\dfrac{\mu^{14}}{(1-\mu^8)^2}TT}{1+\dfrac{\mu^{10}}{(1-\mu^6)^2}TT}\,{.}{.}{.},
\end{aligned}\]}
\noindent und hieraus der allgemeine Ausdruck für $\log\frac tT$ vermittelst einer nach geraden Potenzen von $T$ steigenden Reihe.
\par\bigskip\begin{center}\rule{.18\linewidth}{.4pt}\end{center}
\clearpage
\endgroup
